\section{Introduction}
\label{sec:intro}

Contrastive pretraining learns an image encoder by pulling together two augmented views of the same image
and pushing apart views of different images \citep{chen2020simclr,he2020moco}. It needs no labels. When
labels are available, supervised contrastive learning (\supcon) \citep{khosla2020supcon} uses them to
define more positives: all images of a class are pulled together. With all labels, \supcon is better than
both \simclr and cross-entropy training \citep{khosla2020supcon}.

Most practical problems do not have all labels. They have a small labeled set and a much larger unlabeled
one, and this raises a plain question: \emph{where should the few labels go?} They can go \emph{inside} the
contrastive loss, either by running \supcon on the labeled subset or by mixing labeled and unlabeled images
in one objective \citep{assran2020suncet,inoue2020semi}. Or they can be kept until \emph{after} a
label-free pretraining stage and then used for fine-tuning, as in SimCLR~v2 \citep{chen2020simclrv2}, or
for semi-supervised training \citep{sohn2020fixmatch}.

Prior work gives parts of the answer (\cref{sec:related}), but a controlled comparison has to deal with two
confounds.

\begin{itemize}[leftmargin=1.4em]
\item \textbf{Compute.} A method that trains on the labeled subset and a method that trains on all images do
very different amounts of work if their budgets are counted in passes over their own data: one ``epoch''
over 5,000 images is a tenth of the gradient steps of one epoch over 50,000. Careful studies equalise
iterations \citep{newell2020useful,liu2022self}, but the budget unit is often left unclear: in a small audit
of published comparisons of this kind we could not recover the budget of one arm in 5 of 16 papers, and one
compared arms at equal epochs over pools of different size (\ref{app:audit}).
\item \textbf{Evaluation.} The standard test of a pretrained encoder is a linear probe trained on
\emph{all} training labels \citep{chen2020simclr}. In the label-scarce setting those labels do not exist.
The probe then measures what the features would be worth to someone with an oracle, not what the
practitioner with 5,000 labels can obtain \citep{oliver2018realistic}.
\end{itemize}

This paper removes both confounds. All contrastive methods are written as one loss
(\cref{sec:setup}), and two properties are varied independently in a $2\times2$ design: whether the loss
uses the labels, and whether the pretraining pool is the labeled subset or all images. Every method
receives the same number of gradient steps, from 970 to 97,000, and every encoder is evaluated both with the
all-label probe, which we call the \emph{oracle-label probe}, and with protocols that respect the label
budget. The main experiments use CIFAR-10 with 10\% of the labels (5,000 images), a ResNet-18 and five
seeds; we repeat them on CIFAR-100, on STL-10 with its 100,000 unlabeled images, on Tiny-ImageNet with 200
classes, and with a ResNet-50.

\paragraph{Contributions}
\begin{enumerate}[leftmargin=1.6em]
\item \textbf{A compute-matched $2\times2$ comparison} that separates the use of labels from the choice of
images. It shows that labels inside the loss mainly buy speed. \supcon on the labeled subset is
\R{d:simclr-supcon:lp50:e10}\pt ahead of \simclr after 970 steps, statistically equivalent to it after 9,700
steps, and \R{long:d:simclr-supcon}\pt behind after \R{long:steps} steps
(\cref{sec:results}).
\item \textbf{Evidence that the oracle-label probe overstates label-using encoders.} At 9,700 steps the
oracle-label probe puts a semi-supervised \supcon that trains on all images (\union)
\R{cmp:simclr-union:lp50:e100:diff}\pt ahead of \simclr; fine-tuning on the 5,000 labels puts it
\R{cmp:simclr-union:ft:e100:diff}\pt behind with the default recipe and \R{cmp:simclr-union:fts:e100:diff}\pt behind
with the recipe a practitioner can select, and level with \simclr when each encoder gets its best recipe,
chosen on the test set (\cref{sec:protocol}). With longer training \simclr leads under every probe and fine-tuning recipe, including
that upper bound. Holding out labels before pretraining makes recipe selection possible for the label-using
encoders but does not close the gap, and semi-supervised fine-tuning shrinks all differences to less than one
point.
\item \textbf{A seen/unseen diagnostic.} Comparing the tightness of classes on the labeled images, on the
other training images and on test images shows where the label information ends up. Encoders that use
labels make the labeled images several times tighter than any other images; for \union, the images it trains
on without labels are no tighter than test images (\cref{sec:analysis}).
\end{enumerate}

\paragraph{What the results mean in practice}
Within contrastive pretraining, and in the setting we tested (balanced data, in-distribution or mildly
out-of-distribution unlabeled images, small ResNets), scarce labels should not be spent in a \supcon-style
loss. Pretraining without labels on all images and using the labels afterwards is better, whether
``afterwards'' means fine-tuning, a short \supcon stage (which helps a frozen encoder with a linear head, not
fine-tuning), or semi-supervised training. The last option is
the strongest overall: FixMatch and MixMatch \citep{sohn2020fixmatch,berthelot2019mixmatch} beat every
pretrain-then-fine-tune pipeline in our comparison, and at equal total compute FixMatch from scratch matches
FixMatch started from a contrastive encoder (\cref{sec:where}). We say so plainly,
because the question of this paper is where labels belong, not which contrastive method wins.
